\section{Introduction}

For $n\ge8$, the cycle square $C_n(1,2)$ has vertex set $\Z/n\Z$ and
edges $\{i,i+1\}$ and $\{i,i+2\}$. A signing assigns a value in
$\{\pm1\}$ to each edge. Its signed adjacency matrix $A_\sigma$ is real
symmetric, and we write
\[
 \rho(A_\sigma)=\max\{|\lambda|:\lambda\in\spec(A_\sigma)\},
 \qquad \mu_n=\min_\sigma\rho(A_\sigma).
\]
Both spectral edges matter in this minimization. An upper bound for
$\lambda_{\max}$ alone does not give the required upper bound for $\rho$.

The earlier twisted-optimality conjecture \cite{Suvagiya2026old} compared
$\mu_n$ with the benchmark
\begin{equation}
 \tw(n)^2=4+2\cos\frac{2\pi}{n}+2\cos\frac{4\pi}{n}
 \qquad(n\text{ even},\ n\ge8).
 \label{eq:twisted}
\end{equation}
That preprint has been withdrawn and incorporated into
\cite{Suvagiya2026}. The revised Theorem 26 gives a period-eight signing
with radius
\begin{equation}
 \rstar=\sqrt{4+\sqrt{10+2\sqrt5}}
       =2.7936044933\ldots
 \label{eq:rstar}
\end{equation}
on $8m$ vertices for $m\ge4$, and Conjecture 28 asserts $\mu_{8m}=\rstar$
for $m\ge4$. The construction has positive Hamilton-cycle holonomy;
Remark 27 identifies the opposite holonomy as absent from the reported
search. Changing that holonomy while preserving the triangle signs gives
a strict improvement at every finite size.

Fix the word
\begin{equation}
 t=(1,1,-1,1,-1,-1,1,-1).
 \label{eq:word}
\end{equation}
For $n=8m$, begin with positive step-one edges and step-two signs
$t_{i\bmod8}$. Let $A_m^+$ be this adjacency matrix. Let $A_m^-$ be obtained
by reversing the signs of exactly three edges:
\begin{equation}
 \{n-1,0\},\qquad \{n-2,0\},\qquad \{n-1,1\}.
 \label{eq:seam}
\end{equation}
These are distinct edges for every $n\ge8$.

\begin{theorem}\label{thm:main}
For every integer $m\ge1$, define
\[
 R(s)=\sqrt{4+\sqrt{8+s+\sqrt{26-3s}}}\qquad(-2\le s\le2).
\]
Then
\begin{equation}
 \rho(A_m^+)=R(2)=\rstar,
 \qquad
 \rho(A_m^-)=R\!\left(2\cos\frac{\pi}{m}\right)<\rstar.
 \label{eq:main}
\end{equation}
Moreover, among all signings with the prescribed labeled triangle signs
$t_{i\bmod8}$, the negative-holonomy switching class is the unique
minimizing switching class. In particular,
\[
 \mu_{8m}\le R\!\left(2\cos\frac{\pi}{m}\right)<\rstar
 \qquad(m\ge4),
\]
so Conjecture 28 of \cite{Suvagiya2026} is false as stated.
\end{theorem}

The uniqueness assertion is confined to the prescribed triangle signs.
The theorem supplies no matching lower bound for arbitrary triangle words,
and therefore no classification of unrestricted minimizers. The dispersion relation and exact finite formula are recorded in the
earlier period-eight research manuscript \cite{EarlierPeriodEight}.
The present application concerns the revised conjecture and makes explicit
why its proposed endpoint is missed by every finite negative-holonomy grid.

A separate issue is to obtain a short proof for an explicit family outside
the $8$-divisible subsequence. The following statement is independent of
Theorem~\ref{thm:main}.

\begin{theorem}\label{thm:r2}
Let $k\ge6$, $n=8k+2$, and let $B_n$ be the signed adjacency matrix with
all step-one signs positive and the step-two sign word
\[
 \tau=t^k\mathbin{\|}(1,-1).
\]
Then
\begin{equation}
 \frac{198}{25}I_n-B_n^2\succ0,
 \qquad
 \rho(B_n)^2<\frac{198}{25}<\tw(n)^2.
 \label{eq:r2}
\end{equation}
\end{theorem}

Appendix~\ref{app:r2} gives the full analytic argument and its finite
rational premises. The latter have been checked both through the block
recurrence and independently from the signed graph by scalar Schur
elimination. The theorem is thus a finite-certificate-assisted analytic
result. The signing is inherited from the earlier residue-two construction;
the increment is the analytic closure and its corrected exact certificates.
The proof does not rely on the historical all-even classification.

The character interpretation of twists is part of established spectral
methods. Luo and Roy \cite{LuoRoy2026}, particularly their Theorems 3.6 and
3.13, study graph spectra over homology-character families and endpoint
attainment. Their setting provides context for holonomy, but does not give
an all-signing minimum for this fixed cycle-square support. The
parity-closed-walk identity of Chen, van Dam, and Bu
\cite[Theorem 3.1]{ChenDamBu2024} expresses signed moments averaged over
all signings; an average is likewise not a uniform lower bound for each
signing. Our proofs instead use the explicit fiber polynomial for one
triangle word and exact positive-definiteness certificates for another.
